\section{Related Work}

\paragraph{Theory-of-Mind evaluation in language models.}
Question-answering benchmarks operationalize belief attribution through
controlled stories \citep{le-etal-2019-revisiting}, while conversational
benchmarks test information asymmetry and consistency across related questions
\citep{kim-etal-2023-fantom}. Subsequent stress tests show that high average
accuracy can coexist with failures under trivial or adversarial alterations
\citep{ullman2023trivial,shapira2024cleverhans}. These studies motivate robust
evaluation, but typically retain answer correctness as the principal object.
Our work instead holds story semantics fixed, types the competing answer
sources, and estimates which linguistic description controls arbitration among
them. It therefore addresses controller identification rather than proposing a
broader ToM benchmark.

\paragraph{Discourse accessibility and reference.}
The Givenness Hierarchy links referring form to assumptions about cognitive
status \citep{gundel1993givenness}, and Centering Theory models local discourse
coherence through ranked discourse entities \citep{grosz1995centering}.
Computational coreference work similarly uses mention accessibility to improve
antecedent selection \citep{webster-nothman-2016-using}. Our intervention axes---
description order, explicit names versus ambiguous pronouns, and role labels---
draw on this tradition, but our dependent variable is not human interpretation
or coreference accuracy. We ask how these forms alter a foundation model's
typed answer margin under controlled conflict.

\paragraph{Embedding predicates, factuality, and source commitment.}
Clause-embedding predicates support systematic inferences about the factual
status and source of embedded events. Large-scale lexicosyntactic studies show
that neural models acquire some of these patterns while remaining sensitive to
the lexical and syntactic realization \citep{white-etal-2018-lexicosyntactic}.
Pragmatic evaluations further find that contextual event factuality is not
captured by lexical entailment alone \citep{jiang-de-marneffe-2021-thinks}. We
use reportative and factive predicates as controlled accessibility constraints:
the relevant question is not whether the embedded proposition is globally
entailed, but whether its source remains effective in typed answer arbitration.

\paragraph{Behavioral and compositional evaluation.}
Minimal-pair suites \citep{warstadt-etal-2020-blimp-benchmark}, behavioral test
matrices \citep{ribeiro-etal-2020-beyond}, and contrast sets
\citep{gardner-etal-2020-evaluating} all argue that aggregate test accuracy is
insufficient for characterizing model behavior. Compositional benchmarks test
whether primitive linguistic operations generalize to unseen combinations
\citep{kim-linzen-2020-cogs}. We combine these principles but make the
composition hypothesis quantitative: primitive construction costs are frozen
before unseen compounds are evaluated, and the additive null is rejected by a
preregistered compound-minus-additive inequality. This identifies a behavioral
subadditivity class rather than merely cataloguing another sensitivity.

\paragraph{Language representation as a control interface.}
Recent work argues that language representation design can alter the activation
and organization of model knowledge even without parameter updates
\citep{yang2026schema}. Our results provide a narrow empirical law for that
broader proposal. In controlled belief-binding worlds, representation design
changes source access, but simultaneous linguistic constraints exhibit
saturation rather than independent gains or losses.
